% Aula 07 --- a log-verossimilhança da logística: concavidade, IRLS e separação perfeita.
Beatriz trabalha com marketing e ajusta uma regressão logística para prever se um cliente
responde a uma campanha ($y_i=1$) ou não ($y_i=0$), a partir de $p$ covariáveis. Sejam
$\tilde\x_i=(1,\x_i)\in\R^{p+1}$ o vetor de covariáveis da observação $i$ com um $1$ na
frente, $\tilde{\mathbb X}$ a matriz $n\times(p+1)$ cujas linhas são os $\tilde\x_i^\top$,
$\bbeta=(\beta_0,\beta_1,\dots,\beta_p)$ e
\[
  \pi_i(\bbeta)=\sigma\big(\tilde\x_i^\top\bbeta\big),\qquad
  \sigma(u)=\frac{e^u}{1+e^u},\qquad
  \ell(\bbeta)=\sum_{i=1}^n\Big[y_i\log\pi_i(\bbeta)+(1-y_i)\log\big(1-\pi_i(\bbeta)\big)\Big],
\]
a log-verossimilhança da amostra. O otimizador do pacote que ela usa avisa, numa das bases,
que não convergiu e devolve coeficientes enormes, e ela quer entender por quê.
\begin{itens}
  \item \pts{1,2} Mostre que
        $\ell(\bbeta)=\sum_i\big[y_i\,\tilde\x_i^\top\bbeta-\log\big(1+e^{\tilde\x_i^\top\bbeta}\big)\big]$,
        que o gradiente é $\nabla\ell(\bbeta)=\sum_i\big(y_i-\pi_i\big)\tilde\x_i$ e que a
        hessiana é $-\tilde{\mathbb X}^\top W\tilde{\mathbb X}$, com
        $W=\operatorname{diag}\big(\pi_i(1-\pi_i)\big)$. Conclua que $\ell$ é côncava, que
        todo ponto crítico de $\ell$ é um máximo global e que, se $\tilde{\mathbb X}$
        tem posto $p+1$, há no máximo um maximizador.
  \item \pts{0,6} Mostre que, num maximizador $\widehat\bbeta$, a soma das probabilidades
        ajustadas é igual ao número de clientes que responderam:
        $\sum_i\pi_i(\widehat\bbeta)=\sum_iy_i$. Que leitura isso tem?
  \item \pts{1,0} Não há fórmula fechada para o maximizador, e o otimizador usa o método de
        Newton--Raphson. Supondo $\tilde{\mathbb X}$ de posto $p+1$, escreva o passo de
        Newton que leva $\bbeta$ a $\bbeta^{\mathrm{novo}}$ e mostre que
        $\bbeta^{\mathrm{novo}}$ minimiza a soma de quadrados ponderada
        $\sum_iw_i\big(z_i-\tilde\x_i^\top\bm b\big)^2$ em $\bm b$, com
        $w_i=\pi_i(1-\pi_i)$ e $z_i=\tilde\x_i^\top\bbeta+(y_i-\pi_i)/w_i$, todos
        calculados no $\bbeta$ atual. Por que o método se chama mínimos quadrados
        reponderados iterativamente (IRLS)?
  \item \pts{1,2} Na base problemática existe $\bm b^*\in\R^{p+1}$ com
        $\tilde\x_i^\top\bm b^*>0$ para todo $i$ com $y_i=1$ e $\tilde\x_i^\top\bm b^*<0$
        para todo $i$ com $y_i=0$: um hiperplano separa os que responderam dos que não
        responderam. Mostre que $\ell(\bbeta)<0$ para todo $\bbeta$ e que
        $\ell(t\,\bm b^*)\to0$ quando $t\to\infty$, e conclua que a verossimilhança não
        tem maximizador. Mostre em seguida que, se as duas classes estão presentes, a
        função penalizada $Q(\bbeta)=\ell(\bbeta)-\lambda\sum_{j=1}^p\beta_j^2$, com
        $\lambda>0$ e o intercepto sem penalidade, tem um único maximizador.
\end{itens}

\begin{solucao}
\textbf{(a)} Escreva $u_i=\tilde\x_i^\top\bbeta$. Como $\pi_i=e^{u_i}/(1+e^{u_i})$ e
$1-\pi_i=1/(1+e^{u_i})$, valem $\log\pi_i=u_i-\log(1+e^{u_i})$ e
$\log(1-\pi_i)=-\log(1+e^{u_i})$, e a parcela $i$ de $\ell$ é
\[
  y_i\big[u_i-\log(1+e^{u_i})\big]-(1-y_i)\log(1+e^{u_i})=y_iu_i-\log(1+e^{u_i}),
\]
que é a expressão pedida. A derivada de $u\mapsto\log(1+e^u)$ é $\sigma(u)$, e a de
$\sigma$ é $\sigma(1-\sigma)$. Pela regra da cadeia, com $\partial u_i/\partial\bbeta=\tilde\x_i$,
\[
  \nabla\ell(\bbeta)=\sum_{i=1}^n\big(y_i-\pi_i\big)\tilde\x_i,
  \qquad
  \nabla^2\ell(\bbeta)=-\sum_{i=1}^n\pi_i(1-\pi_i)\,\tilde\x_i\tilde\x_i^\top
   =-\tilde{\mathbb X}^\top W\tilde{\mathbb X}.
\]
Para todo $\bm v\in\R^{p+1}$,
\[
  \bm v^\top\big(\tilde{\mathbb X}^\top W\tilde{\mathbb X}\big)\bm v
  =\sum_{i=1}^n\pi_i(1-\pi_i)\big(\tilde\x_i^\top\bm v\big)^2\ \ge\ 0 ,
\]
e a hessiana é negativa semidefinida em todo ponto: $\ell$ é côncava. Para uma função
côncava e diferenciável, $\ell(\bm b)\le\ell(\bbeta)+\nabla\ell(\bbeta)^\top(\bm b-\bbeta)$
para todos $\bm b,\bbeta$; num ponto crítico o gradiente se anula, e $\ell(\bm b)\le\ell(\bbeta)$
para todo $\bm b$: o ponto crítico é um máximo global. Se $\tilde{\mathbb X}$ tem posto
$p+1$, então $\tilde{\mathbb X}\bm v\ne\bm0$ para $\bm v\ne\bm0$, alguma parcela
$(\tilde\x_i^\top\bm v)^2$ é positiva, e, como $0<\pi_i<1$, a soma acima é positiva: a
hessiana é negativa definida e $\ell$ é estritamente côncava. Se houvesse dois
maximizadores distintos, $\ell$ seria, no ponto médio entre eles, estritamente maior que o
valor máximo, o que é absurdo. Há, portanto, no máximo um.

\smallskip
\textbf{(b)} Num maximizador o gradiente se anula, e a primeira coordenada de
$\nabla\ell$, a que corresponde ao intercepto, é $\sum_i(y_i-\pi_i)\cdot1$. Logo
$\sum_i\pi_i(\widehat\bbeta)=\sum_iy_i$, e, dividindo por $n$, a média das probabilidades
ajustadas é a proporção de clientes que responderam. A logística com intercepto acerta, no
treino, a taxa global de respostas: o que ela pode errar é como essa taxa se distribui entre
os clientes.

\smallskip
\textbf{(c)} O método de Newton para maximizar $\ell$ substitui $\ell$, perto do $\bbeta$
atual, pela sua aproximação de Taylor de segunda ordem e toma o maximizador dela:
$\bbeta^{\mathrm{novo}}=\bbeta-\big[\nabla^2\ell(\bbeta)\big]^{-1}\nabla\ell(\bbeta)$. Com o
item (a), e escrevendo $\bm y$ e $\bm\pi$ para os vetores de $y_i$ e $\pi_i$,
\[
  \bbeta^{\mathrm{novo}}
  =\bbeta+\big(\tilde{\mathbb X}^\top W\tilde{\mathbb X}\big)^{-1}\tilde{\mathbb X}^\top(\bm y-\bm\pi).
\]
A matriz $\tilde{\mathbb X}^\top W\tilde{\mathbb X}$ é invertível porque, pelo item (a),
é positiva definida. Como $W\bm z=W\tilde{\mathbb X}\bbeta+(\bm y-\bm\pi)$, vale
$\tilde{\mathbb X}^\top W\bm z=\tilde{\mathbb X}^\top W\tilde{\mathbb X}\,\bbeta
+\tilde{\mathbb X}^\top(\bm y-\bm\pi)$, e portanto
\[
  \bbeta^{\mathrm{novo}}=\big(\tilde{\mathbb X}^\top W\tilde{\mathbb X}\big)^{-1}\tilde{\mathbb X}^\top W\bm z .
\]
Por outro lado, a soma $S(\bm b)=\sum_iw_i(z_i-\tilde\x_i^\top\bm b)^2$ é uma função
quadrática estritamente convexa de $\bm b$, cujo gradiente,
$-2\tilde{\mathbb X}^\top W(\bm z-\tilde{\mathbb X}\bm b)$, se anula exatamente quando
$\tilde{\mathbb X}^\top W\tilde{\mathbb X}\,\bm b=\tilde{\mathbb X}^\top W\bm z$: são as
equações normais dos mínimos quadrados ponderados, e a solução delas é
$\bbeta^{\mathrm{novo}}$. Cada passo de Newton é, assim, uma regressão de mínimos quadrados
ponderados de uma resposta de trabalho $\bm z$ sobre $\tilde{\mathbb X}$. Os pesos $w_i$ e
a resposta $z_i$ dependem do $\bbeta$ atual e são recalculados a cada passo: daí o nome. Os
pesos são maiores nas observações com $\pi_i$ perto de $1/2$, aquelas sobre as quais o
modelo ainda está mais incerto.

\smallskip
\textbf{(d)} Cada fator da verossimilhança, $\pi_i$ se $y_i=1$ ou $1-\pi_i$ se $y_i=0$, está
em $(0,1)$, porque $\sigma$ toma valores em $(0,1)$. A verossimilhança é um produto de
números em $(0,1)$ e é menor que $1$, e $\ell(\bbeta)<0$ para todo $\bbeta$. Ao longo da
semirreta $t\,\bm b^*$, com $t>0$, cada observação com $y_i=1$ contribui com
$\log\sigma\big(t\,\tilde\x_i^\top\bm b^*\big)$, e $t\,\tilde\x_i^\top\bm b^*\to+\infty$;
cada uma com $y_i=0$ contribui com $\log\big(1-\sigma(t\,\tilde\x_i^\top\bm b^*)\big)$, e
$t\,\tilde\x_i^\top\bm b^*\to-\infty$. Nos dois casos a parcela tende a $\log1=0$, e
$\ell(t\,\bm b^*)\to0$. Assim, $\sup_{\bbeta}\ell(\bbeta)=0$ e esse valor não é atingido em
ponto nenhum: não há maximizador. É o que o otimizador de Beatriz encontrou: cada passo
aumenta $\ell$ afastando $\bbeta$ na direção de $\bm b^*$, e as iterações param pelo limite,
com coeficientes tão grandes quanto o número de passos permitir.

Agora a função penalizada. Ela é contínua. Considere uma sequência $\bbeta^{(m)}$ com
$\norm{\bbeta^{(m)}}\to\infty$, e escreva $\bbeta=(\beta_0,\bbeta_{1:p})$. Passando a uma
subsequência, há dois casos. Se $\norm{\bbeta^{(m)}_{1:p}}\to\infty$, então, como
$\ell\le0$, $Q(\bbeta^{(m)})\le-\lambda\norm{\bbeta^{(m)}_{1:p}}^2\to-\infty$. Se
$\norm{\bbeta^{(m)}_{1:p}}$ fica limitada, então $\abs{\beta_0^{(m)}}\to\infty$. Tome uma
observação $i$ com $y_i=0$ e uma $j$ com $y_j=1$, que existem porque as duas classes estão
presentes; como as parcelas de $\ell$ são negativas,
$Q\le\ell\le\log(1-\pi_i)+\log\pi_j$. Se $\beta_0^{(m)}\to+\infty$, então
$\tilde\x_i^\top\bbeta^{(m)}\to+\infty$ e $\log(1-\pi_i)\to-\infty$; se
$\beta_0^{(m)}\to-\infty$, é $\log\pi_j$ que vai a $-\infty$. Em qualquer caso
$Q(\bbeta^{(m)})\to-\infty$. Logo o conjunto $\{\bbeta: Q(\bbeta)\ge Q(\bm0)\}$ é limitado,
e é fechado porque $Q$ é contínua; numa região compacta, a função contínua $Q$ atinge o
máximo, que é o máximo global, porque fora dela $Q<Q(\bm0)$. A hessiana de $Q$ é
$-\tilde{\mathbb X}^\top W\tilde{\mathbb X}-2\lambda D$, com
$D=\operatorname{diag}(0,1,\dots,1)$. Para $\bm v=(v_0,\bm v_{1:p})\ne\bm0$, se
$\bm v_{1:p}\ne\bm0$ o termo $2\lambda\norm{\bm v_{1:p}}^2$ é positivo; se $\bm v_{1:p}=\bm0$,
então $v_0\ne0$ e
$\bm v^\top\tilde{\mathbb X}^\top W\tilde{\mathbb X}\bm v=v_0^2\sum_i\pi_i(1-\pi_i)>0$. A
hessiana é negativa definida, $Q$ é estritamente côncava, e o maximizador é único, pelo
argumento do item (a). Com separação perfeita, a verossimilhança empurra os coeficientes
para o infinito, e a penalidade os traz de volta: é por isso que a penalização resolve o
problema.
\end{solucao}
